\documentclass[10pt,a4paper]{amsart}
\usepackage[margin=19mm]{geometry}
\usepackage{amsmath,amssymb,mathtools}
\usepackage[scr=rsfs,cal=euler]{mathalfa}
\usepackage{xfrac}
\usepackage[unicode,hidelinks,bookmarks=false]{hyperref}
\newcommand{\ie}{i.e.\ }
\newcommand{\cf}{cf.\ }
\newcommand{\resp}{respectively\ }
\newcommand{\Subsection}{\subsection}
\providecommand{\nobreakdash}{\nobreak\hbox{-}}
\setlength{\emergencystretch}{2em}
\multlinegap0pt
\usepackage{mathtools}
\usepackage{amssymb}
\usepackage{tikz}
\usepackage{bm}
\usepackage{xcolor}
\usepackage{float}
\usepackage{enumerate}
\usepackage[shortlabels]{enumitem}
\newtheorem{lemma}{Lemma}
\newtheorem{theorem}{Theorem}
\newcommand{\R}{\mathbb{R}}
\newcommand{\eps}{\varepsilon}
\newcommand*{\nhid}{n_\textup{hid}}
\newcommand*{\nin}{n_\textup{in}}
\newcommand\norm[1]{\left\lVert#1\right\rVert}
\newcommand\normk[1]{\lVert#1\lVert}
\newcommand*{\nout}{n_\textup{out}}
\title{\textbf{Universal Approximation Constraints of \\ Narrow ResNets: The Tunnel Effect}}
\author{Christian Kuehn, Sara-Viola Kuntz, Tobias W\"ohrer}
\date{Source: arXiv:2603.28591v1 (CC BY 4.0)}
\begin{document}
\maketitle

\noindent\emph{Source and licence.} \textbf{Universal Approximation Constraints of \\ Narrow ResNets: The Tunnel Effect}. Version arXiv:2603.28591v1 (CC BY 4.0), distributed by arXiv under the Creative Commons Attribution 4.0 International licence, http://creativecommons.org/licenses/by/4.0/, as recorded in the arXiv licence metadata for this exact version and shown on the abstract page https://arxiv.org/abs/2603.28591v1. Full licence notice: \texttt{licenses/arxiv-2603.28591-CC-BY-4.0.txt}.\par\smallskip

\noindent\textit{Editorial scope.} Counted unit: the article's theorem th:mlp\_resnet\_error (source0.tex lines 1474--1496). The article's complete original proof of it is delivered with it (source0.tex lines 1498--1526, one \texttt{proof} environment of 29 lines). No mathematics is written, repaired or re-derived: the statement and the proof are the article's own lines, in the article's own notation, and no retained line is substituted or re-flowed.  Retained uncounted as the source-local context the proof needs: lines 1450--1455.  Nothing was omitted from any retained span, so the package carries no omission record.  No retained line contains a citation, so the wrapper carries no bibliography and no bibliography entry.  No retained line contains a figure environment, an image-inclusion command or drawing code, so no image is delivered and the package declares no asset dependency.  Editorial formatting only, in addition: an AMS \texttt{amsart} preamble that restates the article's own declarations and macros used by the retained lines, namely the environments \texttt{lemma}, \texttt{theorem} on separate counters, together with the standard packages those lines need.  The retained environments print in this document as Lemma 1, Theorem 1 in order of appearance, whereas the article prints the same items with the numbering of its own full document.  The complete native source of the article as distributed in the arXiv e-print for this exact version is bundled unchanged as \texttt{sources/197458/source0.tex}.
\begin{lemma}[ResNet Input-Output Map]\label{lem:resnet_inputoutput}
	For a ResNet $\Phi \in \textup{RN}^0_{\eps,\delta}(\mathcal{X},\R^{\nout})$, $\mathcal{X}\subset \R^{\nin}$ open, with $\eps,\delta \geq 0$, for the $l$-th layer it holds that
    \begin{equation}\label{eq:inputoutput}
        h_l = \eps^l \lambda(x) + \delta \sum_{j = 1}^{l} \eps^{l-j}f_j(h_{j-1},\theta_j), \qquad l \in \{1,\ldots,L\}.
    \end{equation}
\end{lemma}

\begin{theorem}[Approximation Error between ResNets and FNNs] \label{th:mlp_resnet_error}   
	Let $0<\eps<1$ and $\delta>0$ and consider a ResNet $\Phi \in \textup{RN}^0_{\eps,\delta}(\mathcal{X},\R^{\nout})$, $\mathcal{X}\subset \R^{\nin}$ open, with update rule
	\begin{equation*}
		h_l = \eps h_{l-1} + \delta f_l(h_{l-1},\theta_l), \qquad l \in \{1,\ldots,L\},
	\end{equation*}
    parameters $\theta_l\in \Theta_l \subset \R^{p_l}$, input transformation $\lambda:\R^{\nin} \rightarrow \R^{\nhid}$ and output transformation $\tilde{\lambda}: \R^{\nhid} \rightarrow \R^{\nout}$.
	Furthermore, consider the corresponding FNN $\overline{\Phi} \in \textup{FNN}^0_{\delta}(\mathcal{X},\R^{\nout})$ with update rule
	\begin{equation*}
		h_l = \delta f_l(h_{l-1},\theta_l), \qquad l \in \{1,\ldots,L\},
	\end{equation*}
	and $h_0 = \lambda(x)$, $h_{L+1} = \tilde{\lambda}(h_L)$. Assume that:
	\begin{itemize}
		\item The residual functions $f_l$ are globally bounded, i.e., there exists a constant $S_f>0$, such that $\norm{f_l(h_{l-1},\theta_l)}_\infty \leq S_f$ for all $h_{l-1} \in \R^{n_{l-1}}$, $\theta_l \in \Theta_l$ and $l \in \{1,\ldots,L\}$. 
        \item The residual functions $f_l$ are globally Lipschitz continuous with Lipschitz constant $K_f$ for all $\theta_l\in\Theta_l$ and $l \in \{1,\ldots,L\}$.  
        \item The input transformation $\lambda$ is globally bounded by the constant $S_\lambda>0$, i.e., $\norm{\lambda}_{\infty,\mathcal{X}}\leq S_\lambda$.
        \item The output transformation $\tilde{\lambda}$ is globally Lipschitz continuous with Lipschitz constant $K_{\tilde{\lambda}}$.
	\end{itemize}
	Then it holds for the global approximation error between the ResNet $\Phi$ and the FNN $\overline{\Phi}$ that
    \begin{equation}\label{eq:error_resnet_mlp}
	   \norm{\Phi-\overline{\Phi}}_{\infty,\mathcal{X}} \leq \eps \cdot K_{\tilde{\lambda}} \left(  (\delta K_f)^{L-1} S_\lambda + (S_\lambda +\delta S_f)  \sum_{j = 0}^{L-2} (\delta K_f)^j\right)  + \mathcal{O}(\eps^2)
    \end{equation}
    as $\eps \rightarrow 0$.
\end{theorem}

\begin{proof}
	Let $x\in \mathcal{X}$ be arbitrary. To estimate $\normk{\Phi-\overline{\Phi}}_{\infty,\mathcal{X}}$, we first bound the distance between the hidden states of the ResNet, denoted by $h_l^{\textup{RN}}$, and the hidden states of the FNN, denoted by $h_l^{\textup{FNN}}$ for $l \in \{1,\ldots,L\}$. The pointwise error at layer $l$ is defined to be $\Delta_l \coloneqq \norm{h_l^{\textup{RN}}-h_l^{\textup{FNN}}}_\infty$. 
    For the first hidden layer it holds 
    \begin{equation*}
        \Delta_1 = \norm{h_1^{\textup{RN}}-h_1^{\textup{FNN}}}_\infty = \norm{\eps \lambda(x) + \delta f_1(\lambda(x),\theta_1) -\delta f_1(\lambda(x),\theta_1)}_\infty = \norm{\eps \lambda(x)}_\infty \leq \eps S_\lambda.
    \end{equation*}
    To determine $\Delta_l$ for $l \in \{2, \ldots, L\}$, we first estimate with Lemma~\ref{lem:resnet_inputoutput} :
    \begin{align*}
        \norm{h_l^{\textup{RN}}}_\infty &= \norm{\eps^l \lambda(x) + \delta \sum_{j = 1}^{l} \eps^{l-j}f_j(h_{j-1}^{\textup{RN}},\theta_j)}_\infty \leq \eps^l \norm{\lambda(x)}_\infty + \delta \sum_{j=1}^{l} \eps^{l-j} \norm{f_j(h_{j-1}^{\textup{RN}},\theta_j)}_\infty \\
        &\leq \eps^l S_\lambda  + \delta S_f  \sum_{j=0}^{l-1} \eps^j 
        \leq S_\lambda  + \frac{\delta S_f}{1-\eps} \eqqcolon H,
    \end{align*}
    where we used the bound of the geometric series for $0 < \eps < 1$. It follows for $l \in \{2,\ldots,L\}$:
    \begin{align*}
        \Delta_l &= \norm{\eps h_{l-1}^{\textup{RN}} + \delta f_l(h_{l-1}^{\textup{RN}}, \theta_l) - \delta f_l(h_{l-1}^{\textup{FNN}}, \theta_l)}_\infty \\
        &\leq \eps \norm{h_{l-1}^{\textup{RN}}}_\infty + \delta \norm{f_l(h_{l-1}^{\textup{RN}}, \theta_l) - f_l(h_{l-1}^{\textup{FNN}}, \theta_l)}_\infty \leq \eps H + \delta K_f \Delta_{l-1},
    \end{align*}
    where we used the Lipschitz continuity of the residual functions $f_l$. As the bound of $\Delta_l$ depends linearly on $\Delta_{l-1}$, we can estimate
    \begin{equation*}
        \Delta_L \leq \eps H + \delta K_f \Delta_{L-1} \leq \ldots \leq (\delta K_f)^{L-1}\Delta_1 + \eps H \sum_{j = 0}^{L-2} (\delta K_f)^j.
    \end{equation*}
    Finally, we use the Lipschitz continuity of the output transformation $\tilde{\lambda}$ to obtain
    \begin{align}
        \norm{\Phi - \overline{\Phi}}_{\infty,\mathcal{X}} &= \sup_{x \in \mathcal{X}}\norm{\tilde{\lambda}(h_L^{\textup{RN}}) - \tilde{\lambda}(h_L^{\textup{FNN}})}_\infty 
        \leq K_{\tilde{\lambda}} \Delta_L \nonumber \\
        &\leq \eps \cdot K_{\tilde{\lambda}}\left( (\delta K_f)^{L-1} S_\lambda +  \left(S_\lambda  + \frac{\delta S_f}{1-\eps}\right) \sum_{j = 0}^{L-2} (\delta K_f)^j\right),  \label{eq:explicit_bound} 
    \end{align}
    such that the statement follows after expanding $\frac{1}{1-\eps}$ in $\eps$.
\end{proof}

\end{document}
